\documentclass[12pt,a4paper]{ctexart}

% ===== 页面设置 =====
\usepackage[top=2.5cm,bottom=2.5cm,left=2.8cm,right=2.8cm]{geometry}
\usepackage{hyperref}
\usepackage{graphicx}
\usepackage{float}
\usepackage{caption}
\usepackage{enumitem}
\usepackage{booktabs}
\usepackage{longtable}
\usepackage{multirow}
\usepackage{array}
\usepackage{xcolor}
\usepackage{colortbl}
\usepackage{listings}
\usepackage{tikz}
\usetikzlibrary{shapes,arrows,positioning,fit}

% ===== 代码高亮 =====
\lstset{
    basicstyle=\ttfamily\small,
    numbers=left,
    numberstyle=\tiny,
    frame=single,
    breaklines=true,
    language=SQL
}

% ===== 超链接 =====
\hypersetup{
    colorlinks=true,
    linkcolor=blue,
    urlcolor=blue
}

\title{\textbf{办公物品管理系统\\数据库设计与实现}}
\author{现代软件实验 + 实验六：综合实验}
\date{\today}

\begin{document}
\maketitle
\tableofcontents
\newpage

% ================================================================
\section{系统概述及需求分析}
% ================================================================

\subsection{业务背景}

某单位办公物品管理存在以下问题：物品采购入库依赖纸质登记，领用审批流程不规范，库存信息无法实时掌握。需要开发一套信息化管理系统，实现办公物品的规范化管理。

\subsection{角色与功能需求}

系统包含两种角色：

\begin{description}
    \item[管理员] 通过桌面应用程序登录，负责用户维护、物品信息管理、入库登记、领用审批和库存查询。
    \item[普通用户] 通过 Web 浏览器登录，只能查看库存和提交领用申请。
\end{description}

\subsection{数据存储要求}

\begin{itemize}
    \item 用户信息存储在文件中（JSON 格式）
    \item 物品、入库、出库等业务数据存储在数据库中（MySQL 8.4.8 LTS）
\end{itemize}

% ================================================================
\section{概念结构设计}
% ================================================================

\subsection{实体识别}

通过对需求的分析，识别出以下实体：

\begin{enumerate}
    \item \textbf{物品类别}（Category）：纸张、文具、刀具、单据、礼品、其它
    \item \textbf{产地}（Origin）：北京、上海、广州等 13 个产地
    \item \textbf{物品}（Item）：办公物品的基本信息和库存数量
    \item \textbf{入库记录}（StockIn）：管理员采购入库的流水记录
    \item \textbf{领用记录}（StockOut）：普通用户申请领用的流水记录
\end{enumerate}

\subsection{实体属性}

\subsubsection{物品类别实体}

\begin{itemize}
    \item 类别 ID（Id）：INT，主键，自增
    \item 类别名称（Name）：VARCHAR(50)，唯一
\end{itemize}

\subsubsection{产地实体}

\begin{itemize}
    \item 产地 ID（Id）：INT，主键，自增
    \item 产地名称（Name）：VARCHAR(100)，唯一
\end{itemize}

\subsubsection{物品实体}

\begin{itemize}
    \item 物品编码（ItemCode）：VARCHAR(20)，主键
    \item 物品名称（ItemName）：VARCHAR(100)，必填
    \item 类别（Category）：INT，外键关联类别表
    \item 产地（Origin）：VARCHAR(100)
    \item 规格（Specification）：VARCHAR(100)
    \item 型号（Model）：VARCHAR(100)
    \item 图片路径（ImagePath）：VARCHAR(500)
    \item 图片数据（ImageData）：LONGBLOB
    \item 库存数量（Quantity）：INT，默认 0
\end{itemize}

\subsubsection{入库记录实体}

\begin{itemize}
    \item 入库流水号（StockInNo）：VARCHAR(20)，主键，格式 SI+日期+序号
    \item 物品编码（ItemCode）：VARCHAR(20)，外键
    \item 购买日期（PurchaseDate）：DATETIME
    \item 入库数量（Quantity）：INT，必须大于 0
    \item 单价（UnitPrice）：DECIMAL(18,2)
    \item 总价（TotalPrice）：DECIMAL(18,2)
\end{itemize}

\subsubsection{领用记录实体}

\begin{itemize}
    \item 领用流水号（StockOutNo）：VARCHAR(20)，主键，格式 SO+日期+序号
    \item 物品编码（ItemCode）：VARCHAR(20)，外键
    \item 领用数量（Quantity）：INT，必须大于 0
    \item 领用人 ID（ApplicantId）：VARCHAR(50)
    \item 申请日期（ApplyDate）：DATETIME
    \item 状态（Status）：INT，0=申请 1=确认 2=驳回
    \item 审批日期（ApproveDate）：DATETIME
    \item 备注（Remark）：VARCHAR(500)
\end{itemize}

\subsection{ER 图}

实体间关系如下：

\begin{itemize}
    \item 一个类别对应多个物品（1:N）
    \item 一个物品对应多条入库记录（1:N）
    \item 一个物品对应多条领用记录（1:N）
    \item 产地与物品之间没有强制外键约束，通过应用程序层面实现取值关联
\end{itemize}

\begin{figure}[H]
\centering
\begin{tikzpicture}[
    entity/.style={rectangle,draw,minimum width=2.4cm,minimum height=3.5cm,align=center},
    relation/.style={diamond,draw,minimum width=1.5cm,minimum height=0.8cm,align=center}
]
    \node[entity] (cat) at (0,0) {\textbf{Categories}\\[2pt] Id (PK)\\ Name (UQ)};
    \node[entity] (item) at (4,0) {\textbf{Items}\\[2pt] ItemCode (PK)\\ ItemName\\ Category (FK)\\ Origin\\ Specification\\ Model\\ ImagePath\\ ImageData\\ Quantity};
    \node[entity] (si) at (8,2) {\textbf{StockIn}\\[2pt] StockInNo (PK)\\ ItemCode (FK)\\ PurchaseDate\\ Quantity\\ UnitPrice\\ TotalPrice};
    \node[entity] (so) at (8,-2) {\textbf{StockOut}\\[2pt] StockOutNo (PK)\\ ItemCode (FK)\\ Quantity\\ ApplicantId\\ ApplyDate\\ Status\\ ApproveDate\\ Remark};
    \node[entity] (ori) at (0,-3) {\textbf{Origins}\\[2pt] Id (PK)\\ Name (UQ)};

    \draw[->,>=stealth] (cat.east) -- node[above] {1:N} (item.west);
    \draw[->,>=stealth] (item.north east) -- node[above,sloped] {1:N} (si.south west);
    \draw[->,>=stealth] (item.south east) -- node[below,sloped] {1:N} (so.north west);
    \draw[->,>=stealth,dashed] (ori.east) -- node[below,sloped] {参考} (item.south west);
\end{tikzpicture}
\caption{办公物品管理系统 ER 图}
\end{figure}

% ================================================================
\section{逻辑结构设计}
% ================================================================

\subsection{关系模式}

将 ER 图转换为关系模式（5 张表）：

\begin{enumerate}
    \item \textbf{Categories}(\underline{Id}, Name)
    \item \textbf{Origins}(\underline{Id}, Name)
    \item \textbf{Items}(\underline{ItemCode}, ItemName, Category, Origin, Specification, Model, ImagePath, ImageData, Quantity)\\
    外键：Category REFERENCES Categories(Id)
    \item \textbf{StockIn}(\underline{StockInNo}, ItemCode, PurchaseDate, Quantity, UnitPrice, TotalPrice)\\
    外键：ItemCode REFERENCES Items(ItemCode)
    \item \textbf{StockOut}(\underline{StockOutNo}, ItemCode, Quantity, ApplicantId, ApplyDate, Status, ApproveDate, Remark)\\
    外键：ItemCode REFERENCES Items(ItemCode)
\end{enumerate}

\subsection{完整性约束}

\subsubsection{实体完整性}

每张表均定义主键约束，确保记录唯一标识：

\begin{itemize}
    \item PK\_Categories：Categories(Id)
    \item PK\_Origins：Origins(Id)
    \item PK\_Items：Items(ItemCode)
    \item PK\_StockIn：StockIn(StockInNo)
    \item PK\_StockOut：StockOut(StockOutNo)
\end{itemize}

\subsubsection{参照完整性}

\begin{itemize}
    \item FK\_Items\_Category：Items.Category → Categories.Id，ON DELETE RESTRICT, ON UPDATE CASCADE
    \item FK\_StockIn\_ItemCode：StockIn.ItemCode → Items.ItemCode，ON DELETE RESTRICT, ON UPDATE CASCADE
    \item FK\_StockOut\_ItemCode：StockOut.ItemCode → Items.ItemCode，ON DELETE RESTRICT, ON UPDATE CASCADE
\end{itemize}

RESTRICT 策略的含义：如果某物品已被入库或领用引用，则不允许删除该物品，必须先删除关联记录。这保证了数据的引用完整性。

\subsubsection{用户定义完整性}

\begin{itemize}
    \item CK\_Items\_Quantity：CHECK (Quantity >= 0)，库存不能为负
    \item CK\_StockIn\_Quantity：CHECK (Quantity > 0)，入库数量必须为正
    \item CK\_StockIn\_UnitPrice：CHECK (UnitPrice >= 0)，单价不能为负
    \item CK\_StockIn\_TotalPrice：CHECK (TotalPrice >= 0)，总价不能为负
    \item CK\_StockOut\_Quantity：CHECK (Quantity > 0)，领用数量必须为正
    \item CK\_StockOut\_Status：CHECK (Status IN (0, 1, 2))，状态只能是 0/1/2
    \item UQ\_Categories\_Name：UNIQUE，类别名称唯一
    \item UQ\_Origins\_Name：UNIQUE，产地名称唯一
\end{itemize}

\subsection{范式分析}

所有表均满足第三范式（3NF）：

\begin{itemize}
    \item 每个属性不可再分（满足 1NF）
    \item 不存在部分函数依赖——所有非主属性完全函数依赖于主键（满足 2NF）
    \item 不存在传递函数依赖——非主属性之间没有依赖关系（满足 3NF）
\end{itemize}

\subsection{约束命名规范}

采用前缀命名规则，提高可读性和可维护性：

\begin{table}[H]
\centering
\begin{tabular}{c|c|c}
\toprule
\textbf{前缀} & \textbf{含义} & \textbf{示例} \\
\midrule
PK\_ & 主键（Primary Key） & PK\_Items \\
FK\_ & 外键（Foreign Key） & FK\_StockIn\_ItemCode \\
CK\_ & CHECK 约束 & CK\_Items\_Quantity \\
UQ\_ & UNIQUE 约束 & UQ\_Categories\_Name \\
IX\_ & 索引（Index） & IX\_StockIn\_ItemCode \\
\bottomrule
\end{tabular}
\caption{约束命名规范}
\end{table}

% ================================================================
\section{物理结构设计}
% ================================================================

\subsection{存储引擎选择}

全部表采用 InnoDB 引擎，理由：

\begin{enumerate}
    \item 支持事务（ACID）—— 入库和领用审批需要事务保证原子性
    \item 支持行级锁 —— 高并发场景下性能更好
    \item 支持外键约束 —— 保证数据引用完整性
    \item 支持崩溃恢复 —— 通过 redo log 保证数据安全
\end{enumerate}

\subsection{索引设计}

\begin{longtable}{c|c|c|c}
\toprule
\textbf{索引名} & \textbf{表} & \textbf{字段} & \textbf{设计理由} \\
\midrule
\endfirsthead
\toprule
\textbf{索引名} & \textbf{表} & \textbf{字段} & \textbf{设计理由} \\
\midrule
\endhead
IX\_Items\_Category & Items & Category & 按类别筛选物品时需要 \\
IX\_Items\_ItemName & Items & ItemName & 按名称模糊搜索物品时需要 \\
IX\_StockIn\_ItemCode & StockIn & ItemCode & 按物品查询入库记录时需要 \\
IX\_StockIn\_PurchaseDate & StockIn & PurchaseDate & 按日期范围查询时需要 \\
IX\_StockOut\_ItemCode & StockOut & ItemCode & 按物品查询领用记录时需要 \\
IX\_StockOut\_ApplicantId & StockOut & ApplicantId & 按用户查询领用记录时需要 \\
IX\_StockOut\_Status & StockOut & Status & 按状态筛选待审批记录时需要 \\
\bottomrule
\caption{索引设计与理由}
\end{longtable}

\subsection{存储估算}

以 100 万条入库记录为例：

\begin{itemize}
    \item 每条记录约 80 字节（流水号 20 + 编码 20 + 日期 8 + 数量 4 + 单价 9 + 总价 9 + 索引开销）
    \item 100 万条 ≈ 80 MB（数据） + 40 MB（索引） ≈ 120 MB
    \item 加上图片数据（ImageData LONGBLOB），每张图片按 200KB 估算
    \item 总磁盘占用：120 MB + 图片数据
\end{itemize}

\subsection{性能优化配置}

\begin{lstlisting}[language=SQL,caption=MySQL 配置建议]
-- InnoDB Buffer Pool：设置为物理内存的 50-70%
SET GLOBAL innodb_buffer_pool_size = 536870912;  -- 512MB

-- 日志缓冲区
SET GLOBAL innodb_log_buffer_size = 16777216;    -- 16MB
\end{lstlisting}

% ================================================================
\section{数据库实施}
% ================================================================

\subsection{建表 DDL}

以核心表 Items 为例：

\begin{lstlisting}[language=SQL,caption=Items 表建表语句]
CREATE TABLE Items (
    ItemCode        VARCHAR(20)     NOT NULL,
    ItemName        VARCHAR(100)    NOT NULL,
    Category        INT             NOT NULL,
    Origin          VARCHAR(100)    NOT NULL DEFAULT '',
    Specification   VARCHAR(100)    NOT NULL DEFAULT '',
    Model           VARCHAR(100)    NOT NULL DEFAULT '',
    ImagePath       VARCHAR(500)    NOT NULL DEFAULT '',
    ImageData       LONGBLOB        NULL,
    Quantity        INT             NOT NULL DEFAULT 0,
    CONSTRAINT PK_Items PRIMARY KEY (ItemCode),
    CONSTRAINT FK_Items_Category FOREIGN KEY (Category)
        REFERENCES Categories(Id)
        ON DELETE RESTRICT ON UPDATE CASCADE,
    CONSTRAINT CK_Items_Quantity CHECK (Quantity >= 0)
) ENGINE=InnoDB DEFAULT CHARSET=utf8mb4;
\end{lstlisting}

\subsection{索引创建}

\begin{lstlisting}[language=SQL,caption=索引创建语句]
CREATE INDEX IX_Items_Category ON Items(Category);
CREATE INDEX IX_Items_ItemName ON Items(ItemName);
CREATE INDEX IX_StockIn_ItemCode ON StockIn(ItemCode);
CREATE INDEX IX_StockIn_PurchaseDate ON StockIn(PurchaseDate);
CREATE INDEX IX_StockOut_ItemCode ON StockOut(ItemCode);
CREATE INDEX IX_StockOut_ApplicantId ON StockOut(ApplicantId);
CREATE INDEX IX_StockOut_Status ON StockOut(Status);
\end{lstlisting}

\subsection{存储过程}

共创建 10 个存储过程，以下是核心的几个：

\begin{lstlisting}[language=SQL,caption=流水号生成存储过程]
CREATE PROCEDURE sp_GenerateStockInNo(OUT outNo VARCHAR(20))
BEGIN
    DECLARE prefix VARCHAR(10);
    DECLARE seq INT;
    SET prefix = CONCAT('SI', DATE_FORMAT(NOW(), '%Y%m%d'));
    SELECT COUNT(*) + 1 INTO seq
    FROM StockIn WHERE StockInNo LIKE CONCAT(prefix, '%');
    SET outNo = CONCAT(prefix, LPAD(seq, 3, '0'));
END
\end{lstlisting}

\begin{lstlisting}[language=SQL,caption=入库事务存储过程]
CREATE PROCEDURE sp_StockIn(
    IN p_ItemCode VARCHAR(20),
    IN p_PurchaseDate DATETIME,
    IN p_Quantity INT,
    IN p_UnitPrice DECIMAL(18,2)
)
BEGIN
    DECLARE v_StockInNo VARCHAR(20);
    DECLARE v_TotalPrice DECIMAL(18,2);
    DECLARE EXIT HANDLER FOR SQLEXCEPTION
    BEGIN
        ROLLBACK;
        RESIGNAL;
    END;

    START TRANSACTION;
        CALL sp_GenerateStockInNo(v_StockInNo);
        SET v_TotalPrice = p_Quantity * p_UnitPrice;

        INSERT INTO StockIn (StockInNo, ItemCode, PurchaseDate,
                             Quantity, UnitPrice, TotalPrice)
        VALUES (v_StockInNo, p_ItemCode, p_PurchaseDate,
                p_Quantity, p_UnitPrice, v_TotalPrice);

        UPDATE Items SET Quantity = Quantity + p_Quantity
        WHERE ItemCode = p_ItemCode;

        SELECT v_StockInNo AS StockInNo;
    COMMIT;
END
\end{lstlisting}

\subsection{触发器}

共创建 4 个触发器，实现库存的自动同步：

\begin{lstlisting}[language=SQL,caption=入库触发器和确认领用触发器]
-- 入库后自动累加库存
CREATE TRIGGER trg_StockIn_AfterInsert
AFTER INSERT ON StockIn
FOR EACH ROW
BEGIN
    UPDATE Items SET Quantity = Quantity + NEW.Quantity
    WHERE ItemCode = NEW.ItemCode;
END;

-- 领用确认后自动扣减库存
CREATE TRIGGER trg_StockOut_AfterUpdate
AFTER UPDATE ON StockOut
FOR EACH ROW
BEGIN
    IF OLD.Status = 0 AND NEW.Status = 1 THEN
        UPDATE Items
        SET Quantity = GREATEST(0, Quantity - NEW.Quantity)
        WHERE ItemCode = NEW.ItemCode;
    END IF;
END;
\end{lstlisting}

\subsection{视图}

\begin{lstlisting}[language=SQL,caption=库存汇总视图]
CREATE VIEW v_StockSummary AS
SELECT
    i.ItemCode, i.ItemName, c.Name AS CategoryName,
    i.Origin, i.Specification, i.Model, i.Quantity,
    CASE
        WHEN i.Quantity = 0 THEN '缺货'
        WHEN i.Quantity < 50 THEN '低库存'
        WHEN i.Quantity < 200 THEN '正常'
        ELSE '充足'
    END AS StockLevel
FROM Items i
JOIN Categories c ON i.Category = c.Id;
\end{lstlisting}

\subsection{种子数据}

预置 6 个类别、13 个产地、10 种物品和示例入库/领用记录，确保系统安装后即可使用。

\subsection{百万条测试数据}

通过存储过程循环生成 100 万条入库记录，用于性能测试。

\begin{lstlisting}[language=SQL,caption=百万条数据生成存储过程（核心循环）]
WHILE i <= 1000000 DO
    SET v_Idx = 1 + FLOOR(RAND() * 10);
    SET v_ItemCode = SUBSTRING_INDEX(
        SUBSTRING_INDEX(v_ItemCodes, ',', v_Idx), ',', -1);
    SET v_Price = 1 + ROUND(RAND() * 499, 2);

    INSERT INTO StockIn (...) VALUES (...);

    IF i % 10000 = 0 THEN COMMIT; END IF;
    SET i = i + 1;
END WHILE;
\end{lstlisting}

% ================================================================
\section{数据库运行和维护}
% ================================================================

\subsection{性能监控}

\begin{lstlisting}[language=SQL,caption=慢查询监控]
-- 查看慢查询配置
SHOW VARIABLES LIKE 'slow_query%';
SHOW VARIABLES LIKE 'long_query_time';

-- 启用慢查询日志
SET GLOBAL slow_query_log = 1;
SET GLOBAL long_query_time = 0.5;  -- 超过 0.5 秒视作慢查询
\end{lstlisting}

\subsection{执行计划分析}

\begin{lstlisting}[language=SQL,caption=EXPLAIN 分析]
EXPLAIN SELECT * FROM StockIn WHERE ItemCode = 'P001';
-- 预期结果：type=ref, key=IX_StockIn_ItemCode, rows≈100,000

EXPLAIN SELECT ItemCode, SUM(TotalPrice) FROM StockIn
GROUP BY ItemCode;
-- 预期结果：type=index, Using temporary
\end{lstlisting}

\subsection{性能测试结果}

在插入 100 万条入库记录后测试：

\begin{table}[H]
\centering
\begin{tabular}{l|c|c}
\toprule
\textbf{测试查询} & \textbf{耗时} & \textbf{说明} \\
\midrule
SELECT COUNT(*) FROM StockIn & $\sim$0.2s & 全表扫描 \\
WHERE ItemCode='P001' & $\sim$0.01s & 走索引 IX\_StockIn\_ItemCode \\
WHERE PurchaseDate BETWEEN '2025-01' AND '2025-06' & $\sim$0.05s & 走索引 IX\_StockIn\_PurchaseDate \\
GROUP BY ItemCode（聚合） & $\sim$0.8s & 全表扫描 + 临时表 \\
JOIN Items 查询 & $\sim$0.3s & 索引 JOIN \\
\bottomrule
\end{tabular}
\caption{百万条记录性能测试结果}
\end{table}

\subsection{备份策略}

\begin{lstlisting}[language=bash,caption=数据库备份命令]
# 全量备份
mysqldump -u root -p --databases OfficeItemsDB > backup.sql

# 仅备份结构
mysqldump -u root -p --no-data OfficeItemsDB > schema.sql

# 定时备份（crontab）
0 2 * * * mysqldump -u root -p密码 OfficeItemsDB > backup_$(date +\%Y\%m\%d).sql
\end{lstlisting}

% ================================================================
\section{总结}
% ================================================================

\subsection{设计要点回顾}

\begin{enumerate}
    \item \textbf{三层完整性保障}：实体完整性（主键）、参照完整性（外键 + RESTRICT）、用户定义完整性（CHECK 约束）
    \item \textbf{库存一致性的双重保障}：应用层事务（EF Core BeginTransaction）+ 数据库层触发器，确保任何方式修改数据都能保持库存正确
    \item \textbf{流水号编码规则}：SI/SO + 日期 + 序号，既保证唯一性又具有业务可读性
    \item \textbf{索引策略}：针对高频查询条件建立索引，百万条数据下索引查询耗时低于 10ms
    \item \textbf{范式设计}：全部表符合 3NF，消除数据冗余和更新异常
\end{enumerate}

\subsection{问题与解决方案}

\begin{description}
    \item[问题一：EF Core 实体跟踪冲突] 单例 DbContext 下，同一主键的实体被 FindAsync 跟踪后，Update 传入新实例会报跟踪冲突。解决方案：GetByCodeAsync 改用 AsNoTracking()，UpdateAsync 前 Detach 已跟踪实体。
    \item[问题二：100 万条数据全量加载卡死] 入库页面 GetAllAsync() 加载全部记录导致界面卡死。解决方案：新增 GetRecentAsync(50) 仅加载最近 50 条，同时库存查询页使用分页。
    \item[问题三：触发器性能影响] 100 万条 INSERT 时每条都触发触发器，耗时显著增加。解决方案：种子数据脚本中临时禁用触发器，插入完成后重新启用。
\end{description}

\subsection{收获与体会}

通过本次数据库设计与实现，我深入理解了：

\begin{itemize}
    \item 从需求分析到物理实现的完整数据库设计流程
    \item ER 图向关系模式的转换方法
    \item 完整性约束（实体、参照、用户定义）的实际应用
    \item 存储过程、触发器、视图、游标等高级数据库对象的使用场景
    \item 索引对查询性能的影响及 EXPLAIN 分析工具的使用
    \item 数据库事务在保证数据一致性中的关键作用
\end{itemize}

\end{document}
